\section{数字护栏、数字锁与研报纠错}
\label{sec:guards}

抽取与计算可以做到确定，但文稿终究要由模型来写。本节介绍把式 \eqref{eq:provenance} 落到实处的几道程序检查：数字护栏（文稿里的数字有没有出处）、指标绑定（数字是否配对到了正确的指标与期间）、外部来源登记、出处分类、数字锁（改写有没有动数字）和内部一致性检查。研报纠错（赛题 5）与护栏共用同一个核查引擎，放在本节最后介绍。

\begin{figure}[t]
  \centering
  \figslot{fig4_guard}{5.0cm}
  \caption{数字护栏与数字锁。左：护栏把文稿中的每个“数字 + 单位”按单位与展示精度换算，到本次运行的结论性产物与外部来源登记中找出处；找不到则列为“无出处数字”，指标或期间配错则列为“绑定不符”，只在来源登记中有出处却未标编号则列为“缺引用”。右：数字锁比较改写前后的数字集合，只允许减少，不允许新增或改动。}
  \label{fig:guard}
\end{figure}

\subsection{数字护栏}
\label{sec:guard}

\paragraph{记号化。}护栏用一个正则表达式从文稿中抽取“数字 + 单位”记号 $x=(r,u)$，单位包括 亿元、万元、元、元/股、元/吨、\%、个百分点、倍 等。年份（19xx、20xx）、一位整数序号、“第 $n$ 页 / 节”这类不带单位的记号不检查。

\paragraph{出处集合。}出处值集合 $V$ \textbf{只取结论性产物}：真值表与分析结果、业绩对比、经营数据（按列单位换算）、DCF 结果、可比公司结果、预测摘要、分部预测、估值计划中的 WACC 与终值参数、股权价值桥各项、以及纠错结论中的正确值。估值产物以百万元计，护栏另按 $\times 10^6$ 换算。两类数值被有意排除在 $V$ 之外：
\begin{itemize}
  \item \textbf{研报草稿中的断言原文。}草稿里的错误数字不能成为出处，否则纠错引擎看过的错数会反过来为文稿背书。
  \item \textbf{中间网格。}三表逐期明细与敏感性矩阵包含成千上万个中间值，把它们纳入 $V$ 会显著提高编造数字“碰巧”命中的概率。
\end{itemize}

\paragraph{按展示精度的容差。}设记号 $x$ 的原文写法有 $d$ 位小数，单位换算因子为 $\phi_u$（亿元为 $10^{-8}$，万元为 $10^{-4}$；百分数额外允许以小数存储的比例，即因子 $\{1, 100\}$），则
\begin{equation}
  \mathrm{pass}(x) \iff \exists\, v\in V,\ \exists\,\phi\in\Phi_u:\ \bigl|\,|v|\cdot\phi - |x|\,\bigr| \le \tfrac12\,10^{-d}\,(1+10^{-9}) .
  \label{eq:guard}
\end{equation}
也就是说，“58.62 亿元”只与真值除以 $10^8$ 后四舍五入到两位小数相比较，“14.49\%”只与百分数值比较。容差随展示精度收紧，是护栏能压低误放行的关键：若把一个编造的数字视为在宽度为 $W$ 的区间内随机取值，而 $V$ 中落在该区间的值有 $n$ 个，则它碰巧通过的概率粗略为 $n\cdot 10^{-d}/W$。精度越高、$V$ 越小，误放行越少；这也是我们宁可把中间网格排除在 $V$ 之外的原因。

\paragraph{指标绑定。}式 \eqref{eq:guard} 只回答“这个数有没有出处”，回答不了“这个数是不是配给了对的指标与期间”——模型完全可能把营业成本的数写成营业收入。护栏因此把文稿当作一份研报，送入纠错引擎（第 \ref{sec:check} 节）按“指标 + 期间”逐条核对。为避免误报，只对口径明确的断言做绑定核对：预测表述、季度与半年度数、跨年序列、分部数据、现金流明细、“由 $a$ 升至 $b$”这类区间表述一律跳过（它们仍要通过式 \eqref{eq:guard} 的出处核对）；分部表述（例如“某产品毛利率”）若句中点名了经营数据表的某一行、且数值等于该行某个已核对的值，则按分部数据核对通过。

\subsection{外部来源登记}
\label{sec:sources}
年报之外的事实只能来自外部来源登记表 \file{research/sources.json}。每条来源记录标题、发布方、网址、发布日期与访问日期，以及若干条原文摘句（不超过 120 字）；摘句中声明的每个数值都必须作为子串出现在摘句里——这与模型抽取“数字须出现在所写页码的原文中”是同一原则。每个检索问题还记录所用检索词与是否找到。护栏把登记表中的数值并入出处集合，但附加一条规则：\textbf{只能在外部来源中找到出处的数字，所在句子必须带对应的 \code{[S编号]}}，否则列为“缺引用”。

\subsection{出处分类：事实、推论与观点}
赛题要求区分事实、推论与观点。我们把这一要求落到每个数字上：护栏按产物给出处值打标签——真值表中按抽取方式分为“报告披露”与“程序计算”，可比公司原始数据与来源登记为“外部来源”，DCF 与预测为“模型预测”，WACC 与终值参数为“估值假设”。文稿中每个有出处的数字归入
\[
  \text{事实}=\{\text{报告披露},\ \text{外部来源}\},\quad \text{推论}=\{\text{程序计算}\},\quad \text{观点}=\{\text{模型预测},\ \text{估值假设}\},
\]
同一数字在多处有出处时优先认作事实。研报附录据此写出全文数字的事实、推论、观点构成。打标签不改变护栏的判定结果；我们在 24 次真实运行上核对过，带标签的出处集合与原出处集合逐值一致。

\subsection{数字锁}
\label{sec:numlock}
复核编辑会话会删句、调序、改措辞。设 $\mathrm{Tok}(\cdot)$ 为去掉注释与来源标记后的“数字 + 单位”记号多重集，$\mathrm{Src}(\cdot)$ 为来源编号集合。数字锁要求
\begin{equation}
  \mathrm{set}\bigl(\mathrm{Tok}(N_{\text{改写}})\bigr)\subseteq \mathrm{set}\bigl(\mathrm{Tok}(N_{\text{初稿}})\bigr)\quad\wedge\quad \mathrm{Src}(N_{\text{改写}})\subseteq \mathrm{Src}(N_{\text{初稿}}).
\end{equation}
数字可以减少，不能新增或改动；来源编号同理。数字锁只看记号本身，数字与指标的配对仍由改写后重跑的护栏核对。需要补数字或改结论的问题，复核编辑不能自己动手，只能写进人工复核清单，或在质检循环中路由回分析与写作侧（第 \ref{sec:refine} 节）。数字锁让我们可以放心地给编辑会话很大的改写自由度：它能改变的只是措辞与结构。

\subsection{内部一致性}
程序只能核对单个数字的出处，核对不了“摘要写买入、正文写卖出”这类前后矛盾。一致性检查从文稿中抽取评级、各情景每股价值、预测年度公司整体的收入、利润、每股收益、毛利率及其增速，\textbf{只有实体、口径、期间、单位、情景都相同时才比较}，两处数值超出展示精度即判为矛盾；外部来源、条件句、敏感性与区间表述、分部口径不进入比较。在七份真实文稿上它零误报，对“摘要评级反转”的人为缺陷全部命中，因此评分标准 v2 把“数字前后矛盾”从评审模型判定改为程序判定。

\subsection{研报纠错（赛题 5）}
\label{sec:check}
纠错时，模型会话先从研报中拆出可核验的数字断言（规则抽取作为补充与交叉检查），每条断言绑定指标、期间、类型、数值与单位。随后按算法 \ref{alg:verify} 判定。

\begin{algorithm}[t]
\caption{断言核查与错误类型诊断（\code{verify\_claim}）}
\label{alg:verify}
\small
\begin{algorithmic}[1]
\Require 断言 $a=(m,p,\tau,x,r,u)$，真值表 $T$
\If{$a$ 为预测性表述} \Return 无法核实 \EndIf
\State $t\gets T.\mathrm{value}(\tau,m,p)$；\textbf{if} $t$ 不存在 \textbf{then return} 无法核实
\State $\epsilon\gets \tfrac12 10^{-d(r)}\phi_u$ \Comment{按展示精度的半个末位}
\If{$|x\phi_u - t|\le\epsilon$} \Return 正确 \EndIf
\If{$\tau=$ 水平值 且 $\exists u'\in\mathrm{Alt}(u): |x\phi_{u'}-t|\le\epsilon'$} \Return 单位错误 \EndIf
\If{$\tau\in\{$同比, 百分点$\}$ 且 $|-x - t|\le\epsilon$} \Return 方向错误 \EndIf
\If{$\exists p'\ne p: |x - T.\mathrm{value}(\tau,m,p')|\le\epsilon$} \Return 期间错配 \EndIf
\If{$\exists m'\in\mathrm{Group}(m)\setminus\{m\}: |x - T.\mathrm{value}(\tau,m',p)|\le\epsilon$} \Return 口径错误 \EndIf
\State \Return 数值错误（附相对偏差与修改建议）
\end{algorithmic}
\end{algorithm}

两个细节体现了“误报控制优先”。其一，错误类型诊断只在\textbf{找到精确匹配的替代解释}时才下结论：断言值换成另一单位、相反方向、另一期间或同组另一指标之后恰好等于真值。找不到时归为“数值错误”，而不是猜一个类型。其二，口径诊断只在容易混用的相关指标组之间进行，例如 \{净利润, 归母净利润, 扣非归母净利润, 利润总额, 营业利润\}、\{营业收入, 营业成本\}、\{毛利率, 净利率\}、\{总资产, 总负债, 所有者权益, 归母权益\}，避免跨指标的巧合匹配造成误判。每条结论都附真值、来源定位（派生值附公式与依赖的事实）和修改建议。
